\documentclass[10pt,a4paper]{amsart}
\usepackage[margin=19mm]{geometry}
\usepackage{amsmath,amssymb,mathtools}
\usepackage[scr=rsfs,cal=euler]{mathalfa}
\usepackage{xfrac}
\usepackage[unicode,hidelinks,bookmarks=false]{hyperref}
\newcommand{\ie}{i.e.\ }
\newcommand{\cf}{cf.\ }
\newcommand{\resp}{respectively\ }
\newcommand{\Subsection}{\subsection}
\providecommand{\nobreakdash}{\nobreak\hbox{-}}
\setlength{\emergencystretch}{2em}
\multlinegap0pt
\usepackage{bm}
\usepackage{bbm}
\usepackage{dsfont}
\usepackage{xspace}
\usepackage{cancel}
\usepackage{mathtools}
\usepackage{amsthm}
\makeatletter
\@ifundefined{theorem}{\newtheorem{theorem}{Theorem}}{}
\makeatother
\title[Basic quantum algebra]{Basic quantum algebra}
\author{Teo Banica}
\date{Source: arXiv:2507.11459v1 (CC BY 4.0)}
\begin{document}
\maketitle

\noindent\emph{Source and licence.} Basic quantum algebra. Version arXiv:2507.11459v1 (CC BY 4.0), distributed under the Creative Commons Attribution 4.0 International licence, as recorded in the arXiv licence metadata for this exact version and shown on the abstract page https://arxiv.org/abs/2507.11459v1. Full rights notice: licenses/arxiv-2507.11459-CC-BY-4.0.txt.\par\smallskip

\noindent\textit{Editorial scope.} Counted unit: the Theorem whose source label is \texttt{None} in the article (source0.tex lines 7702--7706, source label \texttt{None}), delivered together with the article's own complete original proof of it (source0.tex lines 7708--7881, one \texttt{proof} environment of 174 lines). No mathematics is written, repaired or re-derived: statement and proof are the article's own text line for line. Required context retained uncounted: none. Every definition, display and label of those lines is retained unchanged. Omitted from this package and not redistributed: 1--7701, 7707--7707, 7882--17283. Nothing inside a retained excerpt is omitted or altered: every retained body is delivered exactly as the article's lines are written, so no omission receipt of any kind accompanies this package. Citations: the retained excerpts carry no citation command at all, so no bibliography entry of the article is transcribed and no reference is redistributed. Reference adaptation: none was needed and none was made; every reference inside the delivered unit resolves inside it, so no unresolved reference or citation is produced. Printed slips kept literal: no printed word, symbol, hypothesis or number of the counted statement or of its proof was altered. Layout and class: the article's own class is \texttt{amsbook} with packages including amssymb, xy, hyperref; the wrapper is the lane's pinned \texttt{amsart} editorial wrapper at 10pt on A4, which restates the theorem-like environments the retained text needs and adds the article's own macro definitions that the retained text needs (none), with extra wrapper packages (bm, bbm, dsfont, xspace, cancel, mathtools). The delivered numbering is the unit's own. Assets: no retained span contains a figure environment, a graphics-inclusion command, a PDF-page inclusion, a table or a TikZ picture; no image file is copied into this package, the package declares no asset dependency and no image file is redistributed with it. Licence: version arXiv:2507.11459v1 of this article is distributed under the Creative Commons Attribution 4.0 International licence, as recorded in the arXiv licence metadata for this exact version and shown on the abstract page https://arxiv.org/abs/2507.11459v1. A full rights notice is delivered at licenses/arxiv-2507.11459-CC-BY-4.0.txt. Characters: every retained excerpt is pure ASCII, so the delivered engine, XeLaTeX with its default Latin Modern text font, reports no missing character.
\begin{theorem}
We have an isomorphism of compact quantum groups
$$S_4^+=SO_3'$$
given by the Fourier transform over the Klein group $K=\mathbb Z_2\times\mathbb Z_2$.
\end{theorem}

\begin{proof}
This is something which is quite routine, the idea being as follows:

\medskip

(1) Consider the following matrix, coming from the action of $SO_3'$ on $F^4$:
$$u^+=\begin{pmatrix}1&0\\0&u\end{pmatrix}$$

We apply to this matrix the Fourier transform over the Klein group $K=\mathbb Z_2\times\mathbb Z_2$: 
$$v=
\frac{1}{4}
\begin{pmatrix}
1&1&1&1\\
1&-1&-1&1\\
1&-1&1&-1\\
1&1&-1&-1
\end{pmatrix}
\begin{pmatrix}
1&0&0&0\\
0&u_{11}&u_{12}&u_{13}\\
0&u_{21}&u_{22}&u_{23}\\
0&u_{31}&u_{32}&u_{33}
\end{pmatrix}
\begin{pmatrix}
1&1&1&1\\
1&-1&-1&1\\
1&-1&1&-1\\
1&1&-1&-1
\end{pmatrix}$$

By performing the matrix products, we are led to a formula of the following type, with each entry $v_{ij}$ being a certain linear combination of the entries $u_{ij}$:
$$v=\begin{pmatrix}
v_{11}&v_{12}&v_{13}&v_{14}\\
v_{21}&v_{22}&v_{23}&v_{24}\\
v_{31}&v_{32}&v_{33}&v_{34}\\
v_{41}&v_{42}&v_{43}&v_{44}
\end{pmatrix}$$

Our claim now, which will prove the result, is that this matrix is magic, and vice versa, so that the above Fourier transform over the Klein group $K=\mathbb Z_2\times\mathbb Z_2$ converts the relations in Definition 8.5 into the magic relations.

\medskip

(2) In order to prove our claim, let us begin with some preliminaries. For $i,j\in\{1,2,3\}$
with $i\neq j$, we let $<i,j>$ be the unique  element in $\{1,2,3\}$ such that:
$$\{i,j,<i,j>\}=\{1,2,3\}$$

Now let $i,j,k,l \in \{1,2,3\}$ with $i\neq j$ and $k\neq l$. Our claim is that we have:
$$u_{<i,j><k,l>}=u_{ik}u_{jl}+u_{jk}u_{il}$$  

Indeed, for $i,j\in\{1,2,3\}$, let $i_1,i_2,j_1,j_2$ be the elements of $\{1,2,3\}$ satisfying:
$$\{i,i_1,i_2\}=\{j,j_1,j_2\}=\{1,2,3\} $$

Then, the well-known formula for the antipode of the quantum group $SU_3'$ gives:
$$S(u_{ji})=u_{i_1j_1}u_{i_2j_2}+u_{i_2j_1}u_{i_1j_2}$$

But since the matrix $u=(u_{ij})$ is orthogonal, we have the following formula:
$$u_{ij}=S(u_{ij})$$

Thus, we have proved the above claim, regarding the value of $u_{<i,j><k,l>}$.  

\medskip

(3) Let us prove now that $SO_3'$ is the universal quantum group acting on $F^4$. Let $e_1, e_2, e_3, e_4$ be the canonical basis of $F^4$, and consider as well the following basis:
\begin{eqnarray*}
1&=&e_1+e_2+e_3+e_4\\
\varepsilon_1&=&e_1-e_2-e_3+e_4\\
\varepsilon_2&=&e_1-e_2+e_3-e_4\\
\varepsilon_3&=&e_1+e_2-e_3-e_4
\end{eqnarray*}

Observe that this latter basis is precisely the one obtained by using the Fourier transform of the Klein group $K=\mathbb Z_2 \times\mathbb Z_2$, mentioned in (1). We have the following formulae, regarding this latter basis, which define a presentation of the algebra $F^4$:
$$\varepsilon_i^2=1\quad,\quad\forall i$$
$$\varepsilon_i\varepsilon_j=\varepsilon_j \varepsilon_i=\varepsilon_{<i, j>}\quad,\quad \forall i \neq j$$

Consider now the linear map $\alpha :F^4 \to F^4 \otimes F(SO_3')$ defined as follows:
$$\alpha(1)=1\otimes 1\quad,\quad \alpha(\varepsilon_i)=\sum_j\varepsilon_j\otimes u_{ji}$$

Our claim is that $\alpha$ is a coaction. Indeed, it is clear that $\alpha$ is coassociative, so it remains to check that $\alpha$ is a morphism of algebras. First, we have:
\begin{eqnarray*}
\alpha(\varepsilon_i)^2
&=&\sum_{kl}\varepsilon_k\varepsilon_l\otimes u_{ki}u_{li}\\
&=&\sum_k\varepsilon_k^2\otimes u_{ki}^2+\sum_{k\neq l}\varepsilon_k\varepsilon_l\otimes u_{ki}u_{li}\\
&=&1\otimes\left(\sum_ku_{ki}^2\right)+\sum_{k<l}\varepsilon_k\varepsilon_l\otimes(u_{ki}u_{li} +u_{li}u_{ki})\\
&=&1\otimes 1\\
&=&\alpha(\varepsilon_i^2)
\end{eqnarray*}

Also, by using the formula found in (2), we have, for any $i\neq j$:
\begin{eqnarray*}
\alpha(\varepsilon_i)\alpha(\varepsilon_j)
&=&\sum_{kl}\varepsilon_k\varepsilon_l\otimes u_{ki}u_{lj}\\
&=&\sum_k\varepsilon_k^2\otimes u_{ki}u_{kj}+\sum_{k\neq l}\varepsilon_k\varepsilon_l\otimes
 u_{ki}u_{lj}\\
&=&1\otimes\left(\sum_ku_{ki}u_{kj}\right)+\sum_{k<l}\varepsilon_k\varepsilon_l\otimes(u_{ki}u_{lj}+u_{li}u_{kj})\\
&=&\sum_{k<l}\varepsilon_{<k,l>}\otimes u_{<k,l\><i,j>}\\
&=&\sum_k\varepsilon_k\otimes u_{k< i,j>}\\
&=&\alpha(\varepsilon_{<i,j>})
\end{eqnarray*}

We conclude that our map $\alpha$ constructed above is indeed a coaction.

\medskip

(4) Our claim now is that $\alpha$ is the universal coaction on $F^4$. In order to prove this, consider a Hopf algebra $A$ coacting on $F^4$, with coaction denoted as follows:
$$\beta :F^4 \to F^4 \otimes A$$

We set $\varepsilon_0 =1$ and we write, with $x_{j0}=\delta_{j0}$: 
$$\beta(\varepsilon_i)=\sum_j\varepsilon_j\otimes x_{ji}$$

Let $\phi :F^4 \to F$ be the map given by $\phi(e_i) =1/4$, so that $\phi(\varepsilon_i)=0$ if $i>0$. This linear map $\phi$ is $A$-colinear, and so the linear map $F^4 \otimes F^4 \to F$ given by $x\otimes y\to\phi(xy)$ is $A$-colinear too. It follows that the matrix $x = (x_{ij})$ is orthogonal, and so:
$$x_{01}^2+x_{02}^2 + x_{03}^2=0$$

Therefore $x_{0i} = \delta_{0i}$, the matrix $x' =(x_{ij})_{1\leq i,j\leq3}$ is orthogonal and:
$$\beta(\varepsilon_i) = \sum_{j=1}^3 \varepsilon_j \otimes x_{ji}\quad,\quad\forall i\in\{1,2,3\}$$
 
By using these formulae, we have the following computation:
\begin{eqnarray*}
\beta(1)
&=&\beta(\varepsilon_i^2)\\
&=&\beta(\varepsilon_i)^2\\
&=&\sum_{kl}\varepsilon_k\varepsilon_l\otimes x_{ki}x_{li}\\
&=&\sum_k\varepsilon_k^2\otimes x_{ki}^2+\sum_{k\neq l}\varepsilon_k\varepsilon_l\otimes
 x_{ki}x_{li}\\
&=&1\otimes\left(\sum_kx_{ki}^2\right)+\sum_{k<l}\varepsilon_k\varepsilon_l\otimes (x_{ki}x_{li} +x_{li}x_{ki})\\
&=&1\otimes 1+\sum_{k<l}\varepsilon_k\varepsilon_l\otimes(x_{ki}x_{li}+x_{li}x_{ki})
\end{eqnarray*}

But $\beta(1)=1\otimes1$, so we deduce that for $k\neq l$, we have the following formula:
$$x_{ki}x_{li}=-x_{li}x_{ki}$$

By using now the antipode $S$, which satisfies $S(x_{ij})=x_{ji}$ since the matrix $x'$ is orthogonal, we have as well the following formula:
$$x_{ik}x_{il}=-x_{il}x_{ik}$$

(5) Next, for any $i\neq j$, we have the following computation:
\begin{eqnarray*}
\alpha(\varepsilon_{<i,j>})
&=&\alpha(\varepsilon_i \varepsilon_j)\\
&=&\alpha(\varepsilon_i)\alpha(\varepsilon_j)\\
&=&\sum_{kl}\varepsilon_k\varepsilon_l\otimes x_{ki}x_{lj}\\
&=&\sum_k\varepsilon_k^2\otimes x_{ki}x_{kj}+\sum_{k \neq l}\varepsilon_k\varepsilon_l\otimes
 x_{ki}x_{lj}\\
&=&1\otimes\left(\sum_kx_{ki}x_{kj}\right)+\sum_{k<l}\varepsilon_k\varepsilon_l\otimes(x_{ki}x_{lj}+x_{li}x_{kj})\\
&=&\sum_{k<l}\varepsilon_{<k,l>}\otimes(x_{ki}x_{lj}+x_{li}x_{kj})
\end{eqnarray*}

We conclude from this computation that we have the following formula:
$$x_{<k,l><i,j>}=x_{ki}x_{lj}+x_{li}x_{kj}$$

Similarly, since $\varepsilon_i\varepsilon_j=\varepsilon_j\varepsilon_i$, we have as well the following formula:
$$x_{<k,l><i,j>}=x_{kj}x_{li}+x_{lj}x_{ki}$$
 
By combining now these two relations that we found, we get, for $i\neq j$ and $k\neq l$:
$$[x_{ki},x_{lj}]=[x_{kj},x_{li}]$$

Now by using the antipode, we have as well the following formula:
$$[x_{ki},x_{lj}]=[x_{li},x_{kj}]$$

Summarizing, we have established the following equalities:
$$[x_{ki},x_{lj}]=[x_{li},x_{kj}]=[x_{kj},x_{li}]=-[x_{li},x_{kj}]$$

But this shows that we have the following equality:
$$x_{li}x_{kj} = x_{kj}x_{li}$$

(6) We can now finish. Indeed, we have the following computation, in relation with the quantum determinant one condition, from Definition 8.5:
\begin{eqnarray*}
&&x_{11}x_{22}x_{33}+x_{11}x_{23}x_{32}+x_{12}x_{21}x_{33}+x_{12}x_{23}x_{31}+x_{13}x_{22}x_{31}+x_{13}x_{23}x_{31}\\
&=&x_{11}(x_{22}x_{33}+x_{23}x_{32})+x_{12}(x_{21}x_{33}+x_{23}x_{31})+x_{13}(x_{22}x_{31} +x_{23}x_{31})\\
&=&x_{11}^2+x_{12}^2+x_{13}^2\\
&=&1
\end{eqnarray*}

Thus the quantum determinant one condition is satisfied, so we obtain a Hopf algebra morphism 
$F(SO_3') \to A$ commuting with the respective coactions, and we are done.
\end{proof}

\end{document}
